% Generated draft: requires mathematical and source review.
\documentclass[11pt]{article}
\usepackage[margin=25mm]{geometry}
\usepackage{fontspec}
\setmainfont{Latin Modern Roman}
\usepackage{amsmath,amssymb,mathtools,unicode-math}
\setmathfont{Latin Modern Math}
\usepackage{enumitem,needspace,xurl,hyperref}
\setlength{\parindent}{0pt}
\setlength{\parskip}{5pt}
\setlength{\emergencystretch}{3em}
\newcommand{\N}{\mathbb{N}}
\newcommand{\Z}{\mathbb{Z}}
\newcommand{\Q}{\mathbb{Q}}
\newcommand{\R}{\mathbb{R}}
\newcommand{\C}{\mathbb{C}}
\newcommand{\F}{\mathbb{F}}
\newcommand{\A}{\mathbb{A}}
\newcommand{\PP}{\mathbb{P}}
\newcommand{\E}{\mathbb{E}}
\newcommand{\eps}{\varepsilon}
\newcommand{\dd}{\,\mathrm{d}}
\begin{document}
\section*{UnsolvedMath ID: 20000443}
\textbf{Source problem code:} catalogue entry
\subsection{Definitions and mathematical statement}
Colliot-Th\'el\`ene's AIM question asks whether every smooth projective rationally connected variety $X/k$ has a $k$-point when $\operatorname{cd}(k)\leq1$. Rational connectedness is interpreted geometrically. For the characteristic-zero examples below, cohomological dimension at most one means vanishing of Galois cohomology above degree one for torsion discrete modules. The index is
\[
 \operatorname{ind}(X)=\gcd\{[k(x):k]:x\in X\text{ a closed point}\}.
\]
Index one is weaker than existence of a rational point.

The supplied report attributes two negative answers to different constructions over characteristic-zero fields of cohomological dimension one: Ax's smooth quintic hypersurface in $\PP^9$, of index one and without rational points; and a Colliot-Th\'el\`ene--Madore smooth cubic surface of index three and without rational points over their own field. These are not asserted to lie over a common field. The hypersurfaces are in the geometrically rationally connected class in the report's setting.

Its additional fixed-field corollary chooses the single Colliot-Th\'el\`ene--Madore field $F$ and its cubic surface $S$. It claims that every smooth projective geometrically rationally connected $F$-curve has an $F$-point, while for every integer $n\geq2$ there exists a smooth projective geometrically rationally connected $n$-fold over the same $F$ with index three and no $F$-point. The reported construction may be expressed as
\[
 X_n=S\times_F\PP_F^{n-2},\qquad
 \dim X_n=n,\quad\operatorname{ind}(X_n)=3,\quad X_n(F)=\varnothing.
\]
The curve assertion uses the genus-zero/conic interpretation and the vanishing of the Brauer group in cohomological dimension one. This is a dimension dichotomy over one field, rather than a new independent construction of that field.

The source remarks mention a false proof for an intersection of two quadrics in $\PP^4$, a positive finite-field result attributed to Esnault, and a positive curve-function-field result attributed to Harris, Graber and Starr. The latter lacks its algebraically closed constant-field and characteristic hypotheses in the excerpt, so it is not a claim for arbitrary function fields. Further remarks describe hopes that universal torsors over surfaces might be birational to suitable homogeneous spaces, and an earlier cubic-surface point failure with $(\operatorname{Br}k)[3]=0$, attributed to Madore and Colliot-Th\'el\`ene. These hopes are not established implications. The final ``With hindsight'' remark is truncated and gives no additional recoverable assertion.
\subsection{Short English statement}
Ask whether cohomological dimension at most one forces rational points on rationally connected varieties. The supplied literature claims refute this, including a cubic surface of index three. The report propagates that one surface to every dimension at least two over the same field, while genus-zero curves over that field still have rational points.
\subsection{Sources}
\begin{itemize}
\item Local registry: \url{data/unsolvedmath/problems.json}, permanent ID $20000443$, source code \texttt{catalogue entry}; original question, every supplied remark and research summary.
\item AIM workshop ``Rational and integral points on higher dimensional varieties'', Question 11, attributed to Colliot-Th\'el\`ene: \url{https://aimath.org/WWN/qptsurface2/qptsurface2.pdf}. Canonical location \texttt{aim-algebraic-number-theory-notes.json}, \texttt{notes[94]}; original literature field empty. No PDF was opened here.
\item Supplied attributions: Ax for the index-one quintic; Colliot-Th\'el\`ene--Madore for the index-three cubic surface and field; Esnault and Harris--Graber--Starr for positive regimes. No full bibliographic references, field construction or cubic equation are supplied.
\item Supplied machine-generated research attempt 1: counterexample; status refuted at the run, registry solved and classification \texttt{SOLVED-IN-LITERATURE}; rigor metadata ``proved'', novelty confidence low. The non-peer-reviewed fixed-field corollary requires independent verification before citation.
\end{itemize}
\end{document}
